%==============================================================================
%  Параметрическая переформулировка цветового структурного тензора через
%  модель виртуального освещения: эквивалентность и пределы мультипликативного
%  обобщения на RGB-D
%
%  ПОЛНЫЙ РУССКИЙ ТЕКСТ — для внутренней экспертизы (экспертное заключение о
%  возможности открытого опубликования) и последующей подачи в журнал
%  Pattern Recognition and Image Analysis (Pleiades Publishing).
%
%  УДК 004.932
%
%  СТАТУС: полный черновик. Все числа — фактические результаты репозитория
%  (см. раздел «Воспроизводимость»). Перед подачей автору необходимо:
%   1) вычитать и утвердить каждое утверждение (ответственность автора);
%   2) принять решение о соавторстве (см. комментарий в блоке автора);
%   3) проверить библиографические данные источников [9]–[10] (помечены);
%   4) собрать PDF: latexmk -pdf article_ru.tex (локально или в Overleaf).
%
%  Файл использует стандартный класс article и собирается любым базовым
%  дистрибутивом TeX (в среде подготовки TeX отсутствовал — сборка не
%  проверялась; структура проверена программно).
%==============================================================================
\documentclass[11pt,a4paper]{article}

\usepackage[utf8]{inputenc}
\usepackage[T2A]{fontenc}
\usepackage[russian]{babel}
\usepackage[margin=2.5cm]{geometry}
\usepackage{amsmath,amssymb,amsthm}
\usepackage{booktabs}
\usepackage{graphicx}
\usepackage[hidelinks]{hyperref}

% --- теоремные окружения -----------------------------------------------------
\theoremstyle{plain}
\newtheorem{theorem}{Теорема}
\newtheorem{proposition}{Предложение}
\newtheorem{corollary}{Следствие}
\theoremstyle{definition}
\newtheorem{definition}{Определение}
\theoremstyle{remark}
\newtheorem{remark}{Замечание}

% --- обозначения ---------------------------------------------------------------
\newcommand{\Lvec}{\mathbf{L}}
\newcommand{\grad}{\nabla}
\newcommand{\lmax}{\lambda_{\max}}
\newcommand{\lmin}{\lambda_{\min}}
\newcommand{\Mxy}{M_{XY}}
\newcommand{\T}{^{\mathsf T}}

\title{Параметрическая переформулировка цветового структурного тензора
через модель виртуального освещения: эквивалентность и пределы
мультипликативного обобщения на изображения с каналом глубины}

% Блок автора. РЕШЕНИЕ АВТОРА: подтвердить состав и порядок авторов.
% Контактные данные (e-mail) добавляются на этапе подачи по требованиям журнала.
\author{А.\,А.~Панченко\\
  \small Уфимский университет науки и технологий, Уфа, Россия
  % \and К.\,Ф.~Латыпов\\ \small УУНиТ, Уфа, Россия   % <- при соавторстве раскомментировать
}
\date{}

\begin{document}
\maketitle

\begin{abstract}
Выделение границ на цветных изображениях, чувствительное к изолюминантным
переходам, классически решается многоканальным структурным тензором
Ди~Дзензо. В работе анализируется оператор <<виртуального освещения>>,
моделирующий цветовые каналы изображения как скалярные поверхности,
освещаемые с нескольких виртуальных направлений. Показано, что данный оператор не является
самостоятельным: при отключённой канально-асимметричной модуляции его
усреднённый по направлениям отклик есть строго монотонная функция корня из
наибольшего собственного значения тензора Ди~Дзензо, вследствие чего
пороговые бинаризованные карты границ обоих методов совпадают (ранговая
корреляция Спирмена $\rho \approx 0{,}99$ на естественных изображениях и
$1{,}00$ на синтетических границах). Охарактеризован единственный
алгебраически отличный элемент метода --- мультипликативный <<высотный>>
член: доказано, что он структурно неспособен обнаруживать границы,
определяемые только модулирующим каналом, что закрывает расширение на
изображения с каналом глубины в рассматриваемом классе операторов.
Вклад работы --- явная связь между интуитивной физической моделью и
установившимся тензорным формализмом, классификация допустимых обобщений
и открытый воспроизводимый набор эталонных сравнений семи детекторов
границ.

\medskip
\noindent\textbf{Ключевые слова:} выделение границ, цветные изображения,
структурный тензор, градиент Ди~Дзензо, изолюминантные границы, канал
глубины, виртуальное освещение.
\end{abstract}

%==============================================================================
\section{Введение}
\label{sec:intro}

Классические операторы выделения границ --- Собеля, Превитта,
Кэнни~\cite{canny1986} --- работают с полутоновым представлением изображения
и по построению не реагируют на \emph{изолюминантные} границы: переходы
между областями разного цвета, но одинаковой яркости. Систематическое
решение этой проблемы известно с работы Ди~Дзензо~\cite{dizenzo1986},
предложившего рассматривать градиент многоканального изображения ---
например, заданного в цветовом пространстве <<красный--зелёный--синий>>
(RGB) --- через структурный тензор, сумму внешних произведений поканальных
градиентов, и принимать за силу границы корень из его наибольшего
собственного значения.

Наряду с тензорным формализмом существует интуитивно привлекательная
физическая модель: каналы изображения трактуются как скалярные <<поверхности
высот>>, которые освещаются виртуальными источниками с нескольких
направлений; отклик на границу накапливается по направлениям освещения.
Такая модель естественна для физической интуиции и допускает прямые
обобщения: выбор набора направлений, их количества, весов, а также
асимметричную модуляцию отклика третьим каналом. Оператор этого типа был
предложен и реализован автором независимо от тензорного
формализма~\cite{panchenko2026}.

Настоящая работа отвечает на точный вопрос: \emph{является ли оператор
виртуального освещения новым оператором или переформулировкой
существующего?} Ответ нетривиален: эквивалентность не видна из сравнения
формул, поскольку структурный тензор в модели виртуального освещения
никогда не собирается явно --- он присутствует лишь неявно, в виде суммы
внешних произведений внутри скалярного выражения. Мы устанавливаем
эквивалентность строго, характеризуем единственный алгебраически отличный
элемент модели и доказываем предел его применимости, в том числе для
данных с каналом глубины (RGB-D). Независимые повторные
выводы известных результатов --- обычное явление в науке; их научная
ценность возникает тогда, когда точное отношение к первоисточнику
установлено, а границы новых степеней свободы охарактеризованы. Именно это
и составляет предмет данной статьи.

\paragraph{Вклад работы.}
\begin{itemize}
  \item Параметрическая переформулировка цветового структурного тензора, в
        которой направления <<освещения>>, их число и канально-асимметричная
        модуляция выступают как явные, независимо управляемые степени
        свободы (разд.~\ref{sec:decomp}).
  \item Теорема об эквивалентности: при отключённой модуляции оператор
        сводится к магнитуде Ди~Дзензо с точностью до строго монотонного
        преобразования, т.\,е. даёт тождественные пороговые бинаризованные
        карты границ (разд.~\ref{sec:equiv}, теорема~\ref{thm:equiv}).
  \item Теорема о структурном пределе мультипликативного расширения и
        вытекающая из неё классификация допустимых обобщений на данные
        с каналом глубины (разд.~\ref{sec:rgbd}, теорема~\ref{thm:mult}).
  \item Открытый воспроизводимый бенчмарк (синтетические тесты,
        ориентационное исследование, BSDS500, контролируемый RGB-D)
        с полным исходным кодом (разд.~\ref{sec:exp}).
\end{itemize}

%==============================================================================
\section{Структурный тензор Ди~Дзензо}
\label{sec:background}

Для многоканального изображения с каналами $C_k$ Ди~Дзензо~\cite{dizenzo1986}
определяет градиент через структурный тензор
\begin{equation}
\label{eq:dz-tensor}
M \;=\; \sum_{k} (\grad C_k)(\grad C_k)\T
\;=\;
\begin{pmatrix} E & F \\ F & G \end{pmatrix},
\qquad
\begin{aligned}
E &= \textstyle\sum_k (C_{k,x})^2,\\
F &= \textstyle\sum_k C_{k,x} C_{k,y},\\
G &= \textstyle\sum_k (C_{k,y})^2,
\end{aligned}
\end{equation}
принимая за магнитуду границы $\sqrt{\lmax(M)}$, где
$\lmax \ge \lmin \ge 0$ --- собственные значения $M$. Наибольшее собственное
значение отражает максимальное направленное изменение всех каналов
совместно, поэтому хроматические границы без яркостного перепада
обнаруживаются наравне с яркостными. Развитие этого подхода включает
многоспектральные операторы второго порядка Кумани~\cite{cumani1991},
анизотропную диффузию векторнозначных изображений Сапиро и
Рингача~\cite{sapiro1996} и фотометрически инвариантный цветовой тензор ван
де Вейера и др.~\cite{weijer2006}. Геометрическую трактовку изображения как
многообразия, родственную используемой ниже <<поверхностной>> интуиции,
развивает подход Бельтрами~\cite{sochen1998}.

%==============================================================================
\section{Оператор виртуального освещения как параметрическая декомпозиция}
\label{sec:decomp}

Рассмотрим каналы RGB как три скалярные поверхности. Для упорядоченной
тройки каналов $(X, Y, Z)$ и единичного направления <<освещения>>
$\Lvec = (d_x, d_y)\T$ оператор формирует отклик по направлению
\begin{equation}
\label{eq:vl-response}
E_{\Lvec}
\;=\;
\sqrt{(\grad X \cdot \Lvec)^2 + (\grad Y \cdot \Lvec)^2}
\;\cdot\;
\underbrace{\Bigl(1 + \alpha\,\tfrac{C_Z - T}{255}\Bigr)}_{h(C_Z)},
\end{equation}
где $T$ --- глобальный порог (75-й процентиль распределения интенсивностей),
а $\alpha$ управляет канально-асимметричной \emph{высотной} модуляцией
(при $\alpha=0$ имеем $h \equiv 1$). Отклик усредняется по набору из $N$
направлений освещения $\{\Lvec_k\}$ и, опционально, вычисляется для всех
перестановок каналов с последующим слиянием по максимуму.

Свободные параметры конструкции --- набор направлений $\{\Lvec_k\}$, их
число $N$, модуляция $h$ --- образуют явные <<рычаги управления>>. Основной
результат следующего раздела состоит в том, что при $\alpha = 0$ все эти
степени свободы сворачиваются в один скаляр --- наибольшее собственное
значение тензора Ди~Дзензо, --- так что конструкция~\eqref{eq:vl-response}
есть \emph{параметрическая декомпозиция} известного оператора, а не новый
оператор.

%==============================================================================
\section{Эквивалентность тензору Ди~Дзензо}
\label{sec:equiv}

Сведение опирается на одно тождество: для любого $u \in \mathbb R^2$ и
единичного $\Lvec$
\begin{equation}
\label{eq:identity}
(u \cdot \Lvec)^2 \;=\; \Lvec\T\,(u\,u\T)\,\Lvec .
\end{equation}
Суммируя \eqref{eq:identity} по двум градиентным каналам при $\alpha = 0$,
получаем
\begin{equation}
\label{eq:quadform}
(\grad X \cdot \Lvec)^2 + (\grad Y \cdot \Lvec)^2
=
\Lvec\T
\underbrace{\bigl[(\grad X)(\grad X)\T + (\grad Y)(\grad Y)\T\bigr]}_{\Mxy}
\Lvec ,
\end{equation}
где $\Mxy$ --- в точности двухканальный тензор Ди~Дзензо
вида~\eqref{eq:dz-tensor}. Таким образом,
$E_{\Lvec} = \sqrt{\Lvec\T \Mxy\, \Lvec}$ --- направленная компонента
структурного тензора. В базисе собственных векторов $\Mxy$ при
$\Lvec = (\cos\varphi, \sin\varphi)$
\begin{equation}
\label{eq:radial}
E_\varphi \;=\; \sqrt{\lmax \cos^2\varphi + \lmin \sin^2\varphi}.
\end{equation}

\begin{theorem}[об эквивалентности]
\label{thm:equiv}
Пусть $\bar E = \frac{1}{2\pi}\int_0^{2\pi} E_\varphi\, d\varphi$ ---
усреднённый по направлениям отклик~\eqref{eq:radial}. Тогда
\begin{equation}
\label{eq:elliptic}
\bar E \;=\; \frac{2}{\pi}\,\sqrt{\lmax}\;
\mathrm{E}\!\left(\sqrt{1 - \lmin/\lmax}\right),
\end{equation}
где $\mathrm{E}(\cdot)$ --- полный эллиптический интеграл второго рода. При
фиксированном отношении анизотропии $\lmin/\lmax$ величина $\bar E$ строго
возрастает по $\sqrt{\lmax}$. Следовательно, любой процентильный порог
отбирает по $\bar E$ те же пиксели, что и по магнитуде Ди~Дзензо
$\sqrt{\lmax}$; на структурах ранга один ($\lmin = 0$) бинарные карты границ
совпадают тождественно.
\end{theorem}

\begin{proof}
Формула~\eqref{eq:elliptic} получается из~\eqref{eq:radial} подстановкой
$k^2 = 1 - \lmin/\lmax$ и определением $\mathrm E(k) = \int_0^{\pi/2}
\sqrt{1 - k^2 \sin^2\theta}\, d\theta$ с учётом $\pi$-периодичности
подынтегральной функции. Строгая монотонность по $\sqrt{\lmax}$ при
фиксированном $k$ очевидна, поскольку $\mathrm E(k) > 0$. Строго
возрастающее преобразование сохраняет порядок значений пикселей, а значит,
и всякое процентильное пороговое множество. При $\lmin = 0$ имеем
$\bar E = \tfrac{2}{\pi}\sqrt{\lmax}$, т.\,е. прямую пропорциональность
магнитуде Ди~Дзензо.
\end{proof}

\begin{remark}
Дискретное усреднение по $N = 8$ равноотстоящим направлениям, используемое
в практической реализации, является квадратурной аппроксимацией
интеграла~\eqref{eq:elliptic}; ввиду $\pi$-периодичности
$E_\varphi$ восемь узлов дают четыре независимых отсчёта на период, и
погрешность аппроксимации не влияет на ранговый порядок пикселей на
структурах ранга один.
\end{remark}

\begin{proposition}[об избыточности антиподальных направлений]
\label{prop:mode}
Отклик зависит от $\Lvec$ только через квадратичную форму
$\Lvec\T \Mxy\, \Lvec$, инвариантную относительно замены
$\Lvec \mapsto -\Lvec$. Поэтому пополнение набора направлений
$\{\Lvec_1, \Lvec_2\}$ антиподами $\{-\Lvec_1, -\Lvec_2\}$ не меняет
усреднённый отклик: соответствующие конфигурации порождают попиксельно
идентичные карты границ.
\end{proposition}

Это предсказание подтверждено вычислительно: карты для двух и для четырёх
направлений совпадают побитово (см. разд.~\ref{sec:exp}).

\begin{proposition}[о перестановках каналов]
\label{prop:perms}
При $\alpha = 0$ третий канал $Z$ не входит в отклик, и шесть перестановок
$(X, Y, Z)$ сводятся к трём неупорядоченным парам $\{R,G\}$, $\{R,B\}$,
$\{G,B\}$. Слияние по максимуму даёт
$\max\bigl(\bar E_{RG}, \bar E_{RB}, \bar E_{GB}\bigr)$ --- максимум
направленных откликов двухканальных подтензоров полного трёхканального
тензора~\eqref{eq:dz-tensor}.
\end{proposition}

На границах, где доминирует одна пара каналов, максимум по подтензорам
близок к $\sqrt{\lmax}$ полного трёхканального тензора; этим объясняется
эмпирическая ранговая корреляция $\rho \approx 0{,}99$ (а не строго $1$)
с полным оператором Ди~Дзензо на естественных изображениях.

%==============================================================================
\section{Мультипликативное расширение и его предел на изображениях
с каналом глубины}
\label{sec:rgbd}

Единственный алгебраически отличный элемент
конструкции~\eqref{eq:vl-response} --- мультипликативный высотный член
$h(C_Z)$. Если в роли $Z$ выступает четвёртый, семантически выделенный канал
$D$ (глубина в данных RGB-D), можно ожидать, что канально-асимметричная
модуляция даст преимущество перед симметричным тензором. Покажем, что это
не так, причём по структурным причинам.

\begin{theorem}[о структурном пределе мультипликативной модуляции]
\label{thm:mult}
Пусть полный отклик имеет вид
$E_{\Lvec}^{\mathrm{full}} = h(C_Z)\,\sqrt{\Lvec\T \Mxy\, \Lvec}$
с ограниченной функцией $h$. Если в пикселе $\grad X = \grad Y = 0$, то
$E_{\Lvec}^{\mathrm{full}} = 0$ в этом пикселе для любого $\Lvec$,
независимо от значения $C_Z$ и градиента $\grad C_Z$. Следовательно,
оператор не способен реагировать на границу, определяемую только каналом
$Z$ (например, перепад глубины в области однородного цвета).
\end{theorem}

\begin{proof}
При $\grad X = \grad Y = 0$ имеем $\Mxy = 0$, откуда
$\sqrt{\Lvec\T \Mxy\, \Lvec} = 0$ и $h \cdot 0 = 0$ для любой ограниченной
$h$.
\end{proof}

\begin{corollary}[классификация обобщений]
\label{cor:map}
\emph{Аддитивное} включение канала $D$ --- добавление слагаемого
$(\grad D)(\grad D)\T$ в тензор --- по теореме~\ref{thm:equiv} сводится к
многоканальному оператору Ди~Дзензо. Единственное алгебраически отличное,
\emph{мультипликативное} расширение по теореме~\ref{thm:mult} не
обнаруживает границ, определяемых только каналом $D$. Таким образом, в
рассматриваемом классе не существует обобщения, которое одновременно
отличалось бы от многоканального Ди~Дзензо и обнаруживало границы по
глубине; для этого необходимы механизмы, нелинейные по $\grad D$
(например, модуляция вида $h(\|\grad D\|)$ или логическое стробирование
цветового отклика картой глубинных границ).
\end{corollary}

%==============================================================================
\section{Вычислительные эксперименты}
\label{sec:exp}

Все эксперименты воспроизводятся открытым кодом (разд.~<<Воспроизводимость>>).
Метрика качества --- $F_1$ по паросочетанию с допуском на основе
преобразования расстояния. Замеры времени в работе не используются как
аргумент в пользу какого-либо из операторов: как показано в
замечании~\ref{rem:timing}, они определяются деталями реализации и
методикой измерения в большей степени, чем алгоритмической структурой.

\begin{remark}
\label{rem:timing}
Сравнение операторов по времени выполнения требует осторожности, поскольку
доминируют два фактора, не связанных с алгоритмом. Во-первых, реализация
свёртки: вычисление оператора Собеля через общую двумерную корреляцию
примерно вдвое медленнее сепарабельной реализации при идентичном
результате, поэтому <<паритет по скорости>> с оператором виртуального
освещения возникает лишь при сравнении с неоптимальной реализацией
базового оператора. Во-вторых, методика измерения: при поблочной схеме
(все повторы одного оператора, затем все повторы следующего) один и тот же
вызов на одних и тех же данных давал от 8{,}0 до 4{,}0~мс в зависимости от
состояния менеджера памяти и порядка измерения --- разброс, превышающий
большинство сравниваемых различий. Устойчивые оценки требуют чередования
операторов внутри раунда и перемешивания порядка между раундами.
Отдельно отметим, что зависимость времени от числа направлений освещения
$N$ линейна с большой постоянной составляющей: при размере
$256\times256$ измеренная модель имеет вид
$T(N)\approx 8{,}3 + 0{,}47\,N$~мс, то есть при $N=8$ порядка 69\,\%
времени приходится на операции, не зависящие от $N$ (вычисление
градиентов, оценка порога, нормировка), и сокращение числа направлений с
восьми до двух даёт ускорение лишь около $1{,}5$ раза.
\end{remark}

\subsection{Синтетический набор}

Таблица~\ref{tab:synth} приводит $F_1$ (допуск 2 пикселя, изображения
$256\times256$). Операторы Собеля, Превитта и Ди~Дзензо бинаризованы по
95-му процентилю собственной магнитуды; результат на отдельных
синтетических изображениях чувствителен к этому выбору.

\begin{table}[h]
\centering
\caption{Синтетический набор: $F_1$ при допуске 2 пикселя.}
\label{tab:synth}
\begin{tabular}{lcccc}
\toprule
Изображение & Собель & Кэнни & Ди~Дзензо & Виртуальное освещение \\
\midrule
Шахматная доска            & 0{,}000 & 0{,}945 & 0{,}739 & 0{,}933 \\
Концентрические окружности & 1{,}000 & 1{,}000 & 1{,}000 & 1{,}000 \\
Изолюминантные патчи       & 1{,}000 & 0{,}000 & 0{,}909 & 0{,}909 \\
Цветовой круг              & 0{,}510 & 0{,}471 & 0{,}340 & 0{,}887 \\
Размытый диск              & 1{,}000 & 1{,}000 & 1{,}000 & 1{,}000 \\
Полутоновый градиент       & 1{,}000 & 1{,}000 & 0{,}000 & 0{,}000 \\
\midrule
\textbf{Среднее}           & 0{,}752 & 0{,}736 & 0{,}665 & 0{,}788 \\
\bottomrule
\end{tabular}
\end{table}

\begin{remark}[о <<стерильной>> синтетике]
Высокий результат полутонового оператора Собеля на изолюминантных патчах
объясняется не чувствительностью к цвету: остаточная разница яркостей
выбранной пары цветов составляет ${\sim}0{,}3$ градации серого, и после
нормировки карты градиента этот микросигнал растягивается на весь
динамический диапазон, --- но лишь при отсутствии в кадре более сильного
контента. На смешанном контенте (цветовой круг) и на естественных
изображениях эффект исчезает. Отсюда методологический вывод: демонстрация
цветочувствительности корректна на изображениях со смешанным контентом,
а не на однородных парах цветов.
\end{remark}

\subsection{Ориентационное исследование}

На изолюминантных полосах равной яркости при двенадцати ориентациях нормали
(шаг $15^\circ$) оба цветовых оператора обнаруживают границы
($F_1 \approx 0{,}99$ при всех углах), тогда как оператор Кэнни даёт
$F_1 = 0{,}000$ при каждой из двенадцати ориентаций. Конфигурация с
единственным направлением освещения демонстрирует предсказанную <<слепую
зону>>: $F_1$ падает до $0{,}232$ при ориентации $135^\circ$, когда линия
границы параллельна направлению освещения; уже при двух ортогональных
направлениях провал исчезает ($F_1 \ge 0{,}875$ при всех углах), в полном
согласии с формулой~\eqref{eq:radial}. Попиксельное совпадение карт для
двух и четырёх направлений (предложение~\ref{prop:mode}) подтверждено
побитовой проверкой. На <<звезде>> из 16 секторов (все ориентации в одном
кадре) $F_1$ монотонно растёт с числом направлений: $0{,}797$, $0{,}883$,
$0{,}909$ для одного, двух и восьми направлений соответственно; Ди~Дзензо
даёт $0{,}890$, Кэнни --- $0{,}852$.

\subsection{Естественные изображения: BSDS500}

Таблица~\ref{tab:bsds} приводит результаты на полном тестовом наборе
BSDS500~\cite{arbelaez2011} (200 изображений; эталон --- объединение
разметок всех аннотаторов; допуск $0{,}0075 \times$ диагональ; фиксированные
пороги без подбора ODS/OIS, поэтому абсолютные значения ниже публикуемых в
литературе, но сравнение методов между собой проведено в одинаковых
условиях).

\begin{table}[h]
\centering
\caption{BSDS500, тестовый набор (200 изображений), фиксированные пороги.}
\label{tab:bsds}
\begin{tabular}{lccc}
\toprule
Метод & $F_1$ & Точность & Полнота \\
\midrule
Ди~Дзензо                   & \textbf{0{,}575} & \textbf{0{,}654} & 0{,}531 \\
Кэнни                       & 0{,}566 & 0{,}482 & \textbf{0{,}774} \\
Виртуальное освещение       & 0{,}552 & 0{,}630 & 0{,}509 \\
Превитт                     & 0{,}543 & 0{,}582 & 0{,}529 \\
Собель                      & 0{,}541 & 0{,}577 & 0{,}530 \\
\bottomrule
\end{tabular}
\end{table}

Все рассматриваемые рукотворные операторы существенно уступают современным
обучаемым детекторам (на BSDS500 трансформерные архитектуры достигают
ODS ${\approx}0{,}82$--$0{,}84$~\cite{pu2022edter}); работа не претендует
на конкуренцию с ними --- её предмет аналитический.

\subsection{Эмпирическая проверка эквивалентности}

Таблица~\ref{tab:corr} приводит ранговую корреляцию Спирмена между
непрерывными (до бинаризации) картами отклика, подтверждающую
теорему~\ref{thm:equiv}, а также изолирует эффект высотного члена
сравнением $\alpha = 1$ и $\alpha = 0$.

\begin{table}[h]
\centering
\caption{Ранговая корреляция карт отклика (Спирмен).}
\label{tab:corr}
\begin{tabular}{lcc}
\toprule
Тип изображения & ВО$(\alpha{=}0)$ $\sim$ ДД & ВО$(\alpha{=}1)$ $\sim$ ВО$(\alpha{=}0)$ \\
\midrule
Изолюминантный патч                 & 1{,}000 & 1{,}000 \\
Изолюминантные полосы, $30^\circ$   & 0{,}996 & 0{,}996 \\
Естественные фото (BSDS, среднее)   & 0{,}99  & 0{,}97  \\
\bottomrule
\end{tabular}
\end{table}

Высотный член возмущает ранговый порядок лишь на ${\sim}2$--$3\,\%$
пикселей и на естественных изображениях в среднем не улучшает $F_1$
(0{,}552 против 0{,}575 у Ди~Дзензо, табл.~\ref{tab:bsds}).

\subsection{Контролируемый эксперимент с каналом глубины}

Таблица~\ref{tab:rgbd} подтверждает теорему~\ref{thm:mult} на четырёх
контролируемых конфигурациях $256\times256$: оператор виртуального освещения
с глубиной в роли $Z$ даёт $F_1 = 0{,}000$ на чисто глубинной границе, а
результаты при $\alpha = 1$ и $\alpha = 0$ неотличимы во всех случаях;
симметричный четырёхканальный Ди~Дзензо корректно обрабатывает все
конфигурации.

\begin{table}[h]
\centering
\caption{Контролируемый RGB-D: $F_1$ при допуске 2 пикселя. ДД --- Ди~Дзензо,
ВО --- виртуальное освещение с глубиной в роли $Z$.}
\label{tab:rgbd}
\begin{tabular}{lcccc}
\toprule
Конфигурация & ДД-RGB & ДД-RGBD & ВО-RGBD ($\alpha{=}1$) & ВО-RGBD ($\alpha{=}0$) \\
\midrule
Цветовая граница, глубина постоянна & 0{,}769 & 0{,}769 & 0{,}667 & 0{,}667 \\
\textbf{Глубинная граница, цвет постоянен} & 0{,}000 & \textbf{0{,}769} & \textbf{0{,}000} & 0{,}000 \\
Совпадающие цветовая и глубинная    & 0{,}769 & 0{,}769 & 0{,}667 & 0{,}667 \\
Разнесённые цветовая и глубинная    & 0{,}722 & \textbf{0{,}910} & 0{,}636 & 0{,}636 \\
\midrule
\textbf{Среднее}                    & 0{,}565 & \textbf{0{,}790} & 0{,}490 & 0{,}490 \\
\bottomrule
\end{tabular}
\end{table}

%==============================================================================
\section{Связанные работы}
\label{sec:related}

Многоканальные градиенты восходят к Ди~Дзензо~\cite{dizenzo1986};
Кумани~\cite{cumani1991} распространил подход на операторы второго порядка,
Сапиро и Рингач~\cite{sapiro1996} --- на анизотропную диффузию
векторнозначных изображений, ван де Вейер и др.~\cite{weijer2006} ---
на фотометрически инвариантные цветовые тензоры. Геометрическая трактовка
изображения как двумерного многообразия, вложенного в пространство
<<координаты--цвет>>, развита в рамках подхода Бельтрами~\cite{sochen1998}
и концептуально близка используемой здесь модели поверхностей высот.
Современное состояние задачи выделения границ определяется обучаемыми
методами --- от HED~\cite{xie2015} до трансформерных
архитектур~\cite{pu2022edter}; настоящая работа не конкурирует с ними, а
устанавливает точное отношение между двумя рукотворными операторами и
границы их класса. Насколько известно автору, явное сведение модели
виртуального освещения к структурному тензору (теорема~\ref{thm:equiv}) и
характеризация предела мультипликативной модуляции
(теорема~\ref{thm:mult}) в литературе ранее не описывались.

%==============================================================================
\section{Заключение}
\label{sec:conclusion}

Показано, что интуитивный оператор выделения цветовых границ на основе
виртуального освещения при отключённой канально-асимметричной модуляции
является переформулировкой структурного тензора Ди~Дзензо: усреднённый по
направлениям отклик есть строго монотонная функция $\sqrt{\lmax}$
(через полный эллиптический интеграл второго рода), что влечёт совпадение
пороговых бинаризованных карт границ. Единственный алгебраически отличный
элемент --- мультипликативный высотный член --- структурно неспособен
обнаруживать границы, определяемые только модулирующим каналом, что
закрывает расширение на изображения с каналом глубины внутри
рассматриваемого класса и оставляет открытыми лишь нелинейные механизмы. Практическим побочным результатом
стала открытая библиотека из семи детекторов границ с воспроизводимым
бенчмарком.

Более общий вывод: независимый повторный вывод известного оператора
приобретает научную ценность тогда, когда точное отношение к
первоисточнику установлено строго, а возможности и пределы новой
параметризации охарактеризованы теоремами и контролируемыми экспериментами.

%==============================================================================
\section*{Воспроизводимость}

Полный исходный код, генераторы данных и сценарии экспериментов открыты:
\url{https://github.com/Nervni-Sanya/Edge-detection-benchmark}.
Утверждения теоремы~\ref{thm:equiv} воспроизводятся сценарием
\texttt{tests/correspondence\_analysis.py}, теоремы~\ref{thm:mult} ---
\texttt{tests/rgbd\_synthetic\_test.py}, таблицы --- \texttt{tests/testing.py};
полный вывод содержится в \texttt{docs/equivalence.md}.

% --- Обязательные концевые разделы по правилам журнала (порядок фиксирован) --
\section*{Благодарности}

% РЕШЕНИЕ АВТОРА: подтвердить формулировку благодарностей.
Автор выражает благодарность К.\,Ф.~Латыпову за научное руководство и
обсуждение работы. При разработке программной реализации, проведении
вычислительных экспериментов и подготовке рукописи применялись инструменты
генеративного искусственного интеллекта; постановка задачи, доказательства
и выводы проверены автором, который несёт полную ответственность за
содержание работы.

\section*{Источники финансирования}

Данная работа финансировалась за счёт средств бюджета университета.
Никаких дополнительных грантов на проведение или руководство данным
конкретным исследованием получено не было.

\section*{Конфликт интересов}

Автор данной работы заявляет, что у него нет конфликта интересов.

\section*{Соблюдение этических стандартов}

В данной работе отсутствуют исследования человека или животных.

%==============================================================================
% Список литературы упорядочен по алфавиту английской транскрипции фамилии
% первого автора (правило журнала); нумерация в тексте --- автоматическая.
\begin{thebibliography}{10}

\bibitem{arbelaez2011}
P.~Arbel\'aez, M.~Maire, C.~Fowlkes, and J.~Malik,
``Contour detection and hierarchical image segmentation,''
\emph{IEEE Trans. Pattern Anal. Mach. Intell.},
vol.~33, no.~5, pp.~898--916, 2011.
https://doi.org/10.1109/TPAMI.2010.161

\bibitem{canny1986}
J.~Canny,
``A computational approach to edge detection,''
\emph{IEEE Trans. Pattern Anal. Mach. Intell.},
vol.~8, no.~6, pp.~679--698, 1986.
https://doi.org/10.1109/TPAMI.1986.4767851

\bibitem{cumani1991}
A.~Cumani,
``Edge detection in multispectral images,''
\emph{CVGIP: Graphical Models and Image Processing},
vol.~53, no.~1, pp.~40--51, 1991.

\bibitem{dizenzo1986}
S.~Di~Zenzo,
``A note on the gradient of a multi-image,''
\emph{Computer Vision, Graphics, and Image Processing},
vol.~33, no.~1, pp.~116--125, 1986.
https://doi.org/10.1016/0734-189X(86)90223-9

% ПРОВЕРИТЬ ПЕРЕД ПОДАЧЕЙ: точные выходные данные сборника тезисов.
\bibitem{panchenko2026}
А.\,А.~Панченко,
``Метод выделения границ цветных изображений на основе векторного освещения
и адаптивных перестановок каналов,''
в сб. трудов научной конференции ФТИ УУНиТ, Уфа, 2026.

% ПРОВЕРИТЬ ПЕРЕД ПОДАЧЕЙ: страницы и DOI.
\bibitem{pu2022edter}
M.~Pu, Y.~Huang, Y.~Liu, Q.~Guan, and H.~Ling,
``EDTER: Edge detection with transformer,''
in \emph{Proc. IEEE/CVF Conf. Computer Vision and Pattern Recognition (CVPR)},
2022, pp.~1402--1412.

\bibitem{sapiro1996}
G.~Sapiro and D.~L.~Ringach,
``Anisotropic diffusion of multivalued images with applications to color
filtering,''
\emph{IEEE Trans. Image Process.},
vol.~5, no.~11, pp.~1582--1586, 1996.

\bibitem{sochen1998}
N.~Sochen, R.~Kimmel, and R.~Malladi,
``A general framework for low level vision,''
\emph{IEEE Trans. Image Process.},
vol.~7, no.~3, pp.~310--318, 1998.

\bibitem{weijer2006}
J.~van de Weijer, T.~Gevers, and A.~W.~M.~Smeulders,
``Robust photometric invariant features from the color tensor,''
\emph{IEEE Trans. Image Process.},
vol.~15, no.~1, pp.~118--127, 2006.

\bibitem{xie2015}
S.~Xie and Z.~Tu,
``Holistically-nested edge detection,''
in \emph{Proc. IEEE Int. Conf. Computer Vision (ICCV)}, 2015, pp.~1395--1403.

\end{thebibliography}

\section*{Краткие биографии и фотографии авторов}

Панченко Александр Алексеевич (Panchenko Aleksandr Alekseevich) ---
% ЗАПОЛНИТЬ: краткая биография, статус, область научных интересов;
% приложить фотографию. Учёные степени и должности указываются только здесь.
(заполняется автором).

\end{document}
